% fonts/specimen/legibility.tex — decides spec §8 constants
\documentclass{swisstex}
% The class already defines \condensed via \newfontfamily (default
% texgyreheroscn), so re-pointing it at the installed "SwissTeX Grotesk
% Condensed" family requires \renewfontfamily, not \newfontfamily
% (fontspec errors "Command '\condensed' already defined" otherwise —
% see fonts/specimen/specimen.tex for the same finding, confirmed
% empirically during Task 8).
\setmainfont{SwissTeX Grotesk}
\renewfontfamily\condensed{SwissTeX Grotesk Condensed}
\setrunningtitle{Legibility gate}
\begin{document}
\swisstitle{Prüfdruck}{Lesbarkeitsprobe}{Fussnoten- und Schmalsatz-Entscheid}
\section{Fussnotensatz: 8 auf 12 gegen 8 auf 10{,}125}
{\condensed\fontsize{8}{12}\selectfont
Die Wissenschaft hat festgestellt, dass die Trennlinie auf der Textachse
steht und die Fussnote in der Schmalschrift läuft. Zwölf Punkt Durchschuss
ist das bisherige Mass; es bindet erst nach acht Rasterzeilen wieder an.\par}
\gridskip{1}
{\condensed\fontsize{8}{10.125}\selectfont
Die Wissenschaft hat festgestellt, dass die Trennlinie auf der Textachse
steht und die Fussnote in der Schmalschrift läuft. Zehneinviertel Punkt
Durchschuss ist das starke Verhältnis drei zu vier; es bindet nach drei
Rasterzeilen wieder an.\par}
\section{Schmalsatz bei 0{,}71-Kompression}
\marg{Randglosse in der Schmalschrift bei voller Kompression: bleibt die
Type bei 7{,}5 Punkt offen genug?}
Der Grundtext dient als Umfeld. Zu prüfen ist die Marginalie links: wirken
die Binnenräume der komprimierten Schnitte bei Glossengrösse noch offen,
oder verlangt der Schmalsatz eine mildere Kompression durch Interpolation
innerhalb der eigenen Schnitte?
\colophon{Entscheid bitte im Spec §8/§12.2 nachtragen: Fussnotenmass und
Kompressionsfrage.}
\end{document}
